\FloatBarrier
\section{Physical Origins of Transverse Spin Driving}
\label{app:transverse-drive-origins}
\label{app:electric-spin-driving}

Section~\ref{sec:transverse-spin-control} specifies a transverse
term in the effective spin Hamiltonian. Its physical origin is
contained in the off-diagonal matrix element of the applied
perturbation, after any required orbital or valley reduction.
In the static spin basis, write this contribution as
\begin{equation}
 H_\perp(t)=d_x(t)\sigma_x+d_y(t)\sigma_y,\qquad
 d_x(t)-i d_y(t)
 =\bra{\uparrow}\delta H_{\mathrm{eff}}(t)\ket{\downarrow}.
 \label{eq:transverse-drive-matrix-element}
\end{equation}
Here $d_x$ and $d_y$ have energy units. The drive used in the
main text corresponds to $d_x(t)=\hbar\Omega\cos(\omega t+\phi)$ and
$d_y(t)=0$. The microscopic control determines these coefficients;
the rotating-frame calculation in Appendix~\ref{app:spin-driving}
then determines the resulting spin dynamics.

\subsection{Direct magnetic-field driving}
\label{sec:drive-direct-magnetic}

An oscillating magnetic field couples directly to the electron's
magnetic moment. For a positive isotropic effective $g$ factor
$g_{\mathrm e}$, equal to $g$ in the main text, take
\[
 \mathbf B(t)=B_0\hat{\mathbf z}
       +B_1\cos(\omega t+\phi)\hat{\mathbf x}.
\]
The Zeeman interaction gives
\begin{equation}
 \begin{aligned}
 H(t)&=\frac{E_Z}{2}\sigma_z
       +\frac{g_{\mathrm e}\mu_B B_1}{2}
          \cos(\omega t+\phi)\sigma_x,\\
 E_Z&=g_{\mathrm e}\mu_B B_0,\qquad
 \Omega_{\mathrm{mag}}=\frac{g_{\mathrm e}\mu_B|B_1|}{2\hbar}.
 \end{aligned}
 \label{eq:direct-magnetic-spin-drive}
\end{equation}
The transverse part has a nonzero matrix element between the
Zeeman states. Comparing its coefficient with
Eq.~\eqref{eq:esr-lab} gives the resonant Rabi frequency
$\Omega_{\mathrm{mag}}$ within the RWA. A field component parallel
to the static quantization axis instead changes the splitting.
The spin-only description also requires the isolation conditions
of Section~\ref{sec:platform-encoding}.

\subsection{Electric driving in a magnetic-field gradient}
\label{sec:drive-magnetic-gradient}

An electric field acts directly on the electron's position. For
charge $-e$, with $e>0$, its interaction is
$H_{\mathrm{el}}=e\mathcal E(t)x\otimes I_{\mathrm{spin}}$.
If spin--orbit coupling is neglected, confinement is spin
independent and the Zeeman field is uniform, the two spin states
share an orbital--valley state
$\ket{\chi_0}$. Their electric-field matrix element factorizes as
\begin{equation}
 \bra{\chi_0,\uparrow}H_{\mathrm{el}}\ket{\chi_0,\downarrow}
 =e\mathcal E(t)\bra{\chi_0}x\ket{\chi_0}
   \langle\uparrow|\downarrow\rangle=0.
 \label{eq:electric-selection-rule}
\end{equation}
Thus orbital motion alone does not couple these two spin states.
In a transverse magnetic-field gradient, however, a displaced
electron samples a different transverse field. The same electric
motion then supplies the time-dependent spin coupling. A
micromagnet can produce this gradient~\cite{PioroLadriere2006,Kawakami2014}.
Spin transitions driven through electrically induced motion are
called electric-dipole spin resonance (EDSR).

Take the undriven dot centre as $x=0$, and approximate the local
field by $\mathbf B(x)\simeq B_0\hat{\mathbf z}
+b_\perp x\hat{\mathbf x}$, where
$b_\perp=(\partial B_x/\partial x)_0$.
For motion along $x$, use harmonic confinement with effective
mass $m^*$, orbital frequency $\omega_0$ and momentum $p_x$.
An electric field $\mathcal E(t)=\mathcal E_{\mathrm{ac}}\cos\omega t$
then gives the combined orbital and spin Hamiltonian
\begin{equation}
 H=\frac{p_x^2}{2m^*}+\frac{m^*\omega_0^2x^2}{2}
   +e\mathcal E_{\mathrm{ac}}x\cos(\omega t)
   +\frac{E_Z}{2}\sigma_z
   +\frac{g_{\mathrm e}\mu_B b_\perp x}{2}\sigma_x.
 \label{eq:micromagnet-model}
\end{equation}
The last term makes the direction of the spin splitting depend
on position. We treat its mixing of orbital states as weak:
$g_{\mathrm e}\mu_B|b_\perp|\ell_0\ll\hbar\omega_0$, with
ground-state position width
$\ell_0=\sqrt{\hbar/(2m^*\omega_0)}$.
The driven displacement must also stay in the region where both
the harmonic potential and the constant gradient are valid.

To leading order in this weak mixing, the orbital displacement
can be calculated without spin backaction. Ehrenfest's equation
for $x_c(t)=\langle x\rangle_t$ is a forced-oscillator equation:
\begin{equation}
 \begin{aligned}
 m^*\ddot x_c+m^*\omega_0^2x_c
    &=-e\mathcal E_{\mathrm{ac}}\cos(\omega t),\\
 x_c(t)&=x_{\mathrm{ac}}\cos(\omega t),\qquad
 x_{\mathrm{ac}}=-\frac{e\mathcal E_{\mathrm{ac}}}
                         {m^*(\omega_0^2-\omega^2)}.
 \end{aligned}
 \label{eq:driven-displacement}
\end{equation}
The solution shown omits free orbital oscillations, as appropriate
to a drive switched on smoothly and far from orbital resonance.
For $\omega\ll\omega_0$, it reduces to
$x_{\mathrm{ac}}\simeq-e\mathcal E_{\mathrm{ac}}/(m^*\omega_0^2)$:
the orbital approximately follows the moving potential minimum.

The moving orbital samples the transverse field $b_\perp x_c(t)$.
Projecting the spin-dependent terms of
Eq.~\eqref{eq:micromagnet-model} onto that orbital therefore
replaces $x$ by $x_c(t)$. The resulting spin Hamiltonian and the
Rabi frequency obtained by comparison with
Eq.~\eqref{eq:calibration-electric-spin-drive} are
\begin{equation}
 H_{\mathrm{spin}}(t)\simeq\frac{E_Z}{2}\sigma_z
  +\frac{g_{\mathrm e}\mu_B b_\perp x_{\mathrm{ac}}}{2}
        \cos(\omega t)\sigma_x,\qquad
 \Omega_{\mathrm{mm}}
   =\frac{g_{\mathrm e}\mu_B|b_\perp x_{\mathrm{ac}}|}{2\hbar}.
 \label{eq:micromagnet-rabi}
\end{equation}
The effective drive thus has the same form as a directly applied
oscillating magnetic field. Its transverse direction matters:
a purely longitudinal gradient would modulate the Zeeman
splitting without producing this off-diagonal term.

\subsection{Electric driving through spin--orbit coupling}
\label{sec:drive-spin-orbit}

Spin--orbit coupling connects orbital motion to spin even when
the applied Zeeman field is uniform. Consider a one-dimensional
model in which momentum couples to a transverse spin component
through $\alpha p_x\sigma_y$. The coefficient $\alpha$ has units
of velocity, and the driven Hamiltonian is
\begin{equation}
 H=H_{\mathrm{orb}}+e\mathcal E(t)x
      +\frac{E_Z}{2}\sigma_z+\alpha p_x\sigma_y,\qquad
 H_{\mathrm{orb}}=\frac{p_x^2}{2m^*}+\frac{m^*\omega_0^2x^2}{2}.
 \label{eq:so-model}
\end{equation}
The electric field produces the orbital displacement already
calculated in Eq.~\eqref{eq:driven-displacement}. To extract its
spin-dependent effect, it is useful to remove the momentum
coupling to first order by a change of basis~\cite{Golovach2006}.

For nonzero $\alpha$, introduce the signed spin--orbit length
$\ell_{\mathrm{so}}=\hbar/(m^*\alpha)$ and the anti-Hermitian
generator $A=ix\sigma_y/\ell_{\mathrm{so}}$.
Expanding the unitary transformation gives
$\widetilde H=e^AHe^{-A}=H+[A,H]+\cdots$.
The required commutators are
\begin{equation}
 \begin{aligned}
 [A,H_{\mathrm{orb}}]&=-\alpha p_x\sigma_y,\qquad
 [A,e\mathcal E(t)x]=0,\\
 \left[A,\frac{E_Z}{2}\sigma_z\right]
     &=-\frac{E_Zx}{\ell_{\mathrm{so}}}\sigma_x.
 \end{aligned}
 \label{eq:so-commutators}
\end{equation}
The first cancels the original momentum coupling. The electric
potential is unchanged, while the Zeeman term acquires a
position-dependent transverse component. To first order, the
transformed Hamiltonian is consequently
\begin{equation}
 \widetilde H=H_{\mathrm{orb}}+e\mathcal E(t)x
      +\frac{E_Z}{2}\sigma_z
      -\frac{E_Zx}{\ell_{\mathrm{so}}}\sigma_x+O(\alpha^2).
 \label{eq:so-transformed}
\end{equation}
This form makes the effect of electric motion explicit: the
transverse coupling now depends on the electron's displacement.

For $E_Z,\hbar\omega\ll\hbar\omega_0$ and
$\ell_0,|x_{\mathrm{ac}}|\ll|\ell_{\mathrm{so}}|$, we may project
onto the slowly moving orbital and set $x=x_c(t)$ in the spin
interaction. In the transformed spin basis, this yields
\begin{equation}
 \widetilde H_{\mathrm{spin}}(t)\simeq\frac{E_Z}{2}\sigma_z
  -\frac{E_Z x_{\mathrm{ac}}}{\ell_{\mathrm{so}}}
       \cos(\omega t)\sigma_x,\qquad
 \Omega_{\mathrm{so}}=\frac{E_Z|x_{\mathrm{ac}}|}
                        {\hbar|\ell_{\mathrm{so}}|}.
 \label{eq:so-rabi}
\end{equation}
The retained spin states include the small orbital admixtures
implicit in the basis transformation. In this model the drive
vanishes when $\alpha=0$ or $E_Z=0$; interface, valley and
anisotropic spin--orbit effects can modify the model in silicon.

Equations~\eqref{eq:micromagnet-rabi} and~\eqref{eq:so-rabi}
identify the transverse coefficients produced by the two electric
mechanisms. Once the orbital reduction is valid, both can be
inserted into the same spin dynamics derived in
Appendix~\ref{app:spin-driving}. Within the RWA, their
constant-envelope spin-flip probabilities are given by
Eq.~\eqref{eq:rabi-probability}, with $\Omega$ replaced by
$\Omega_{\mathrm{mm}}$ or $\Omega_{\mathrm{so}}$.

\paragraph{Gate-dependent Zeeman coupling.}
The one-dimensional spin--orbit model is one approximation to the
electrical response. More generally, orbital, valley and spin--orbit
effects can be incorporated in a gate-dependent effective
$g$ matrix~\cite{Venitucci2018GMatrix}. In a fixed reference basis
for the retained doublet, write the Zeeman interaction to first
order in the static applied field $\mathbf B_0$ as
\begin{equation}
 H_Z(V)=\mathbf d(V)\cdot\boldsymbol\sigma,\qquad
 \mathbf d(V)=\frac{\mu_B}{2}\mathbf g(V)\mathbf B_0,\qquad
 \mathbf n=\frac{\mathbf d(V_0)}{|\mathbf d(V_0)|}.
 \label{eq:gate-dependent-zeeman}
\end{equation}
The static splitting is $E_Z=2|\mathbf d(V_0)|$.
For a small modulation $V(t)=V_0+v_{\mathrm{ac}}\cos\omega t$,
linearization gives an oscillating term
$\delta\mathbf d\cdot\boldsymbol\sigma\cos\omega t$.
Its transverse component and resonant Rabi frequency are
\begin{equation}
 \begin{aligned}
 \delta\mathbf d&=v_{\mathrm{ac}}
        \left.\frac{\partial\mathbf d}{\partial V}\right|_{V_0},
 &\delta\mathbf d_\perp
   &=\delta\mathbf d-\mathbf n(\mathbf n\cdot\delta\mathbf d),\\
 \Omega_g&=\frac{|\delta\mathbf d_\perp|}{\hbar}.
 \end{aligned}
 \label{eq:g-matrix-transverse-drive}
\end{equation}
This result assumes the RWA and a drive sufficiently weak and
far below transitions to excluded orbital and valley states for
the doublet reduction to remain valid. A scalar $g(V)$ changes only
the longitudinal splitting; transverse driving requires a change
in the direction of $\mathbf d$. The $g$ matrix describes the
projected microscopic response. A spin--orbit contribution already
included in it must therefore be counted only once.

\subsection{Programming the drive amplitude and phase}
\label{sec:drive-phase-programming}

For the electric-motion mechanisms, the displacement entering the spin
Hamiltonian must be related to the voltage waveform at the dot.
Separate the voltage that sets the operating point from the
oscillating component used for spin rotation:
\begin{equation}
 \begin{aligned}
 V_g(t)&=V_g^{\mathrm{DC}}+V_g^{\mathrm p}(t)+v_d(t),\\
 v_d(t)&=a(t)\cos(\omega t+\phi).
 \end{aligned}
 \label{eq:measurement-gate-waveform}
\end{equation}
Here $V_g^{\mathrm{DC}}$ is the constant bias and
$V_g^{\mathrm p}(t)$ is a slowly varying pulse used for loading,
holding and readout. During a rotation, their sum fixes the
holding point. The drive envelope $a(t)$ denotes the voltage
that actually reaches the dot, including the effects of cables
and filters on the programmed waveform.

At that operating point, let $F_v$ be the force per unit drive
voltage, so that $F_vv_d(t)=-e\mathcal E(t)$ in linear response.
Substituting this force into Eq.~\eqref{eq:driven-displacement},
for a carrier slow compared with the orbital frequency, gives
\[
 x_c(t)\simeq\chi_x a(t)\cos(\omega t+\phi),
 \qquad \chi_x=\frac{F_v}{m^*\omega_0^2}.
\]
The envelope must also vary slowly enough for the orbital to
follow. For the magnetic-gradient mechanism, this displacement
produces
\begin{equation}
 B_{x,\mathrm{ac}}(t)=B_1(t)\cos(\omega t+\phi),\qquad
 B_1(t)=b_\perp\chi_xa(t),
 \label{eq:calibration-transverse-field-from-voltage}
\end{equation}
and hence the signed drive coefficient
\[
 \Omega(t)=\frac{g_{\mathrm e}\mu_Bb_\perp\chi_xa(t)}{2\hbar}.
\]
This relation places the device voltage directly in
Eq.~\eqref{eq:calibration-electric-spin-drive}. The sign of
$b_\perp\chi_x$ can be absorbed into the phase convention. For
the spin--orbit mechanism, the same displacement is inserted into
Eq.~\eqref{eq:so-rabi}.

The phase controls a different part of the rotation: pulses with
the same frequency and envelope can act about different transverse
axes. To set it, a controller combines two versions of the carrier
separated by a quarter period. Their independently chosen,
slowly varying amplitudes give
\begin{equation}
 v_d(t)=v_I(t)\cos\omega t-v_Q(t)\sin\omega t,\qquad
 v_I(t)=a(t)\cos\phi,\quad
 v_Q(t)=a(t)\sin\phi.
 \label{eq:calibration-iq-phase-control}
\end{equation}
Substituting $v_I$ and $v_Q$ yields
$v_d(t)=a(t)\cos(\omega t+\phi)$, the waveform in
Eq.~\eqref{eq:measurement-gate-waveform}. Changing their ratio
and signs therefore sets $\phi$ at fixed carrier frequency and,
through Eq.~\eqref{eq:calibration-single-spin-hamiltonian}, sets
the rotation axis. An IQ mixer driven by an arbitrary-waveform
generator implements this combination. A direct digital
synthesizer can generate the same waveform from a programmed
phase word. No separate static gate voltage is needed to set
the phase.

The phase of each pulse is defined relative to a common carrier
clock. For pulse $k$, the controller turns on an envelope $a_k(t)$
with offset $\phi_k$; turning that envelope off does not reset
the clock. The difference $\phi_2-\phi_1$ therefore remains set
even if the drive is zero between the pulses. A fixed phase shift
from the cables, filters and mixer adds the same offset to the
device phases. It can be absorbed into the definition of $X$,
leaving the relative rotation axes unchanged.

The relation between programmed relative phase and rotation axis can be
checked with two calibrated resonant $\pi/2$ pulses. Prepare
$\ket0$, give the first pulse phase zero and the second phase
$\varphi$, and neglect detuning and decoherence during the
sequence. Equation~\eqref{eq:calibration-single-spin-hamiltonian}
then gives the rotating-frame operators
$U_0=e^{-i\pi X/4}$ and
$U_\varphi=e^{-i\pi(\cos\varphi X-\sin\varphi Y)/4}$.
Their product has excited-state amplitude
\[
 \bra1U_\varphi U_0\ket0
 =-\frac{i}{2}(1+e^{-i\varphi}),
\]
so the ideal measured probability is
\begin{equation}
 P_1(\varphi)=\frac{1+\cos\varphi}{2}.
 \label{eq:calibration-phase-offset-scan}
\end{equation}
The maximum at $\varphi=0$ corresponds to the same axis for both
pulses, and the minimum at $\varphi=\pi$ to opposite axes;
$\varphi=\pm\pi/2$ selects orthogonal axes. With pulse angles and
detuning controlled, a shifted or distorted fringe indicates an
error in the relation between programmed and device phase, such
as an IQ quadrature error. Correcting that relation makes the
programmed phase offsets produce the intended gate axes.
